\documentclass[11pt]{amsart}
\usepackage[T1]{fontenc}
\usepackage{lmodern}
\usepackage{amsmath,amssymb,amsthm,mathtools}
\usepackage[margin=1.1in]{geometry}
\usepackage{booktabs,array,enumitem,xcolor}
\usepackage[expansion=false]{microtype}
\usepackage[pdfusetitle,colorlinks=true,linkcolor=blue!50!black,
  citecolor=blue!50!black,urlcolor=blue!50!black]{hyperref}

\newtheorem{theorem}{Theorem}[section]
\newtheorem{proposition}[theorem]{Proposition}
\newtheorem{lemma}[theorem]{Lemma}
\newtheorem{corollary}[theorem]{Corollary}
\theoremstyle{definition}
\newtheorem{definition}[theorem]{Definition}
\newtheorem{example}[theorem]{Example}
\theoremstyle{remark}
\newtheorem{remark}[theorem]{Remark}

\title{A General Mathematical Article}
\author{Author Name}
\date{\today}

\begin{document}
\begin{abstract}
State the question, the main result, its hypotheses, and the limits of the result.
\end{abstract}
\maketitle
\setcounter{tocdepth}{1}
\tableofcontents

\section{Introduction}\label{sec:introduction}
Introduce the central question and the result of Section~\ref{sec:main}. A sample citation is \cite{lamport}.

\section{Main result}\label{sec:main}
\begin{definition}
An object is a structure satisfying the hypotheses specified in the statement below.
\end{definition}

\begin{theorem}
Under the stated hypotheses, the construction has the claimed property.
\end{theorem}

\begin{proof}
Replace this proof by the complete argument, including all quantifiers and assumptions.
\end{proof}

\section{Conclusion}\label{sec:conclusion}
Record the scope, examples, and remaining questions without enlarging the theorem.

\bibliographystyle{plain}
\bibliography{references}
\end{document}
